\documentclass[11pt]{article}
\usepackage[T1]{fontenc}
\usepackage[margin=1in]{geometry}
\usepackage{amsmath}
\usepackage{hyperref}
\hypersetup{colorlinks=true,linkcolor=blue,citecolor=blue,urlcolor=blue}

\title{A Short Survey of Line-Breaking Algorithms}
\author{M. Alvarez and J. Chen}
\date{}

\begin{document}
\maketitle

\begin{abstract}
We survey three classic approaches to paragraph line breaking and note how
each interacts with incremental re-typesetting. Every source is cited forward
in the text before the bibliography is generated by BibTeX.
\end{abstract}

\tableofcontents

\section{Introduction}
\label{sec:intro}

The greedy line breaker fills each line as far as possible before moving on,
while the optimal-fit algorithm of Knuth and Plass~\cite{plass81} minimizes
total badness across the whole paragraph. Later work recast the same problem
for incremental editors~\cite{adapton}, a view we return to in
Section~\ref{sec:incremental} on page~\pageref{sec:incremental}.

\section{Optimal fit}
\label{sec:optimal}

The details of optimal fit appear in \emph{The \TeX book}~\cite{knuth84}.
Closing the loop on hyphenation required a separate exception
dictionary~\cite{hyphen}, while the dissertation behind the algorithm records
the early experiments~\cite{plass-thesis}.
\nocite{*}

\section{Incremental re-breaking}
\label{sec:incremental}

As previewed in Section~\ref{sec:intro}, incremental algorithms reuse the
breaks computed for unchanged text~\cite{adapton}. Their correctness rests on
the same badness model as optimal fit~\cite{plass81}, since the badness of a
line depends only on its own contents:
\begin{equation}
  \label{eq:badness}
  b = 100 \left\lvert \frac{t}{s} \right\rvert^3,
\end{equation}
where $t$ is the shortfall and $s$ the available stretch.

\bibliographystyle{plain}
\bibliography{refs}
\end{document}
